\documentclass[10pt,a4paper]{article}

% ---------- Packages ----------
\usepackage[margin=0.55in]{geometry}
\usepackage{titlesec}
\usepackage{enumitem}
\usepackage[hidelinks]{hyperref}
\usepackage{xcolor}
\usepackage{tabularx}
\usepackage{array}
\usepackage{needspace}

% ---------- Page Setup ----------
\pagestyle{empty}
\setlength{\parindent}{0pt}
\setlength{\parskip}{0pt}
\renewcommand{\baselinestretch}{0.96}
\urlstyle{same}

% ---------- Colors ----------
\definecolor{sectiongray}{HTML}{222222}
\definecolor{textgray}{HTML}{333333}

% ---------- Section Formatting ----------
\titleformat{\section}
  {\large\bfseries\color{sectiongray}}
  {}
  {0em}
  {}
  [\vspace{-4pt}\titlerule]

\titlespacing*{\section}{0pt}{8pt}{5pt}

% ---------- List Formatting ----------
\setlist[itemize]{
    leftmargin=1.2em,
    itemsep=2pt,
    topsep=2pt,
    parsep=0pt,
    partopsep=0pt
}

% ---------- Custom Commands ----------
\newcommand{\resumeSubheading}[4]{
    \vspace{3pt}
    \textbf{#1} \hfill #2\\
    \textit{#3} \hfill \textit{#4}
    \vspace{2pt}
}

\newcommand{\projectHeading}[3]{
    \Needspace{7\baselineskip}
    \textbf{#1} \hfill \href{#2}{GitHub}\\[-1pt]
    \textit{#3}
    \vspace{-2pt}
}

\newcommand{\educationHeading}[4]{
    \textbf{#1} \hfill #2\\
    #3 \hfill \textit{#4}
    \vspace{2pt}
}

% ---------- Document ----------
\begin{document}

% ---------- Header ----------
\begin{center}
    {\LARGE \textbf{VEERABABU SUTAPALLI}}\\[6pt]
    {\normalsize \textbf{Machine Learning \& GenAI Engineer}}\\[3pt]
    Hyderabad, India \textbar{}
    \href{mailto:isutapalliveerababu@gmail.com}{isutapalliveerababu@gmail.com} \textbar{}
    \href{https://veerababu-sutapalli-portfolio.vercel.app/}{Portfolio}\\[2pt]
    \href{https://www.linkedin.com/in/veera-sutapalli}{linkedin.com/in/veera-sutapalli} \textbar{}
    \href{https://github.com/veera491}{github.com/veera491} \textbar{}
    \href{https://www.hackerrank.com/svb491}{hackerrank.com/svb491} \textbar{}
    \href{https://leetcode.com/u/veera491/}{leetcode.com/u/veera491}
\end{center}

\vspace{-6pt}

% ---------- Professional Summary ----------
\section{Professional Summary}
Computer Science postgraduate focused on machine learning, GenAI, LLM systems, and practical AI application development. Experienced in distributed LLM inference research, AI/LLM product-integration support, large-scale data processing, and ML/NLP projects using Python, PyTorch, Hugging Face, Scikit-learn, and PySpark. T-shaped engineer with supporting strengths across data engineering, backend APIs, full-stack systems, and cloud-backed delivery.

% ---------- Technical Skills ----------
\section{Technical Skills}
\begin{tabularx}{\textwidth}{>{\bfseries}p{3.6cm} X}
AI \& Machine Learning & Python, pandas, NumPy, Scikit-learn, PyTorch Basics, Supervised Learning, Feature Engineering, Cross Validation, Model Evaluation, Explainable AI \\
GenAI \& NLP & Hugging Face Transformers, LLM Inference, NLP, Text Classification, Prompt Engineering, Embeddings, RAG Basics, Structured Outputs, LIME \\
Data Engineering & SQL, MySQL, PySpark, HDFS, ETL, Distributed Processing, Data Cleaning, Parquet \\
Backend \& Full-Stack & Flask, Django, React, Next.js, JavaScript, TypeScript, REST APIs, Firebase \\
Cloud, Testing \& Tools & Azure VMs, AWS, Docker, Git, GitHub, Postman, Playwright, Selenium, Jupyter/Colab \\
\end{tabularx}

% ---------- Experience ----------
\section{Experience}

\resumeSubheading
{Nordic Stone Studio}
{Malm{\"o}, Sweden}
{Programming Intern}
{Feb 2024 -- Jun 2024}
\begin{itemize}
    \item Contributed to Tile Cozy's editor and product workflows by quickly understanding an existing Python, JavaScript, Django, and React codebase and implementing application features.
    \item Supported LLM and Generative AI integration into product systems by preparing data inputs, contributing implementation support, and prototyping AI-assisted workflow components.
    \item Collaborated through team meetings, technical discussions, and Git/GitHub workflows while consistently completing assigned development tasks within internship timelines.
\end{itemize}

\resumeSubheading
{JNTU Incubation Centre, Innovation Research Centre}
{Kakinada, India}
{Big Data and Machine Learning Engineer Intern}
{Jan 2022 -- Jul 2022}
\begin{itemize}
    \item Reduced computation time by 30\% while processing 100M+ healthcare and insurance records by engineering PySpark-based distributed workflows over HDFS.
    \item Improved data consistency and reduced project time by 20\% by building ETL pipelines for patient, hospital, and insurance-claim analytics.
    \item Supported insurance-eligibility prediction and web-based result delivery by preparing ML-ready datasets and collaborating on Python/MySQL-backed implementation.
\end{itemize}

\resumeSubheading
{JNTU Incubation Centre, Innovation Research Centre}
{Kakinada, India}
{Mobile Application Developer Intern}
{Jun 2021 -- Dec 2021}
\begin{itemize}
    \item Developed Python/Flask REST API and MySQL-backed workflows for the Agricultural Advisory Board application, supporting category-wise digitization of farmer grievances.
    \item Reduced analysis time by 40\% through MySQL query optimization and improved reliability using unit checks, Selenium automation, and Postman API validation.
\end{itemize}

% ---------- Projects ----------
\section{Selected Projects}

\projectHeading
{AI-Generated Review \& Comment Detection}
{https://github.com/veera491/ai-generated-review-detection}
{Python, NLP, Random Forest, Pseudo-Labelling, Explainable AI}
\begin{itemize}
    \item Built text-classification workflows across 487,898 labelled and 443,079 unlabelled records using cleaning, normalization, lemmatization, tokenization, EDA, and pseudo-labelling experiments.
    \item Evaluated baseline, pseudo-labelled, and tuned Random Forest models using accuracy, precision, recall, F1-score, and AUC-ROC; achieved 0.85 accuracy and 0.94 AUC-ROC after tuning.
    \item Improved model transparency through LIME, feature-importance analysis, and partial dependence plots.
\end{itemize}

\projectHeading
{Decision Support System for Port Efficiency Analysis}
{https://github.com/veera491/port-efficiency-decision-support-system}
{Python, Flask, pandas, Data Analytics, ML Regression, Visualization}
\begin{itemize}
    \item Contributed to a Flask-based decision-support system analysing 1,330,169 maritime-operation records across 25 columns, including transit, waiting, dwell, vessel-speed, and source-destination patterns.
    \item Compared Linear Regression, Decision Tree Regressor, and Gradient Boosting Regressor using MAE, MSE, and R-squared, and translated findings into stakeholder-facing analytics views.
\end{itemize}

\projectHeading
{Android Malware Detection using Machine Learning}
{https://github.com/veera491/android-malware-detection-ml}
{Python, Scikit-learn, SVM, Random Forest, KNN, Cybersecurity}
\begin{itemize}
    \item Compared Decision Tree, Random Forest, SVM, and KNN classifiers for Android malware detection using accuracy, precision, recall, F1-score, AUC-ROC, and statistical significance testing.
    \item Identified SVM as the strongest overall model, reaching 97.11\% accuracy and leading F1-score and AUC-ROC performance, while Random Forest achieved the highest precision.
\end{itemize}

\projectHeading
{Cloud-Based Brain Stroke Detection}
{https://github.com/veera491/cloud-brain-stroke-detection-flask-aws}
{Python, Flask, Docker, AWS ECR/ECS, Cloud Monitoring}
\begin{itemize}
    \item Developed and containerized a Flask-based brain-stroke detection application, pushed the Docker image to AWS ECR, and deployed it through AWS ECS.
    \item Monitored CPU and memory utilisation through AWS metrics and documented the end-to-end cloud deployment workflow.
\end{itemize}

% ---------- Research & Publications ----------
\section{Research \& Publications}

\noindent
\textbf{Scaling Distributed Inference of BLOOMZ-560M Lightweight Large Language Models using Petals Tool}
\hfill \href{https://www.diva-portal.org/smash/record.jsf?pid=diva2%3A2052191}{Publication}\\
\textit{Master's Thesis, Blekinge Institute of Technology, 2026}
\begin{itemize}
    \item Analysed distributed LLM inference across 1--10 Azure VM nodes using 42,989 prompts and sequential, pipelined, and microbatched execution strategies.
    \item Identified pipelined execution as the strongest scaling strategy with approximately 1.6$\times$ mean speedup at 10 nodes, while observing a 13$\times$ increase in network footprint.
\end{itemize}

\noindent
\textbf{Predicting Risk Level in Life Insurance Application}
\hfill \href{https://www.diva-portal.org/smash/record.jsf?pid=diva2%3A1784294}{Publication}\\
\textit{Bachelor's Thesis, Blekinge Institute of Technology, 2023}
\begin{itemize}
    \item Compared Logistic Regression, Decision Tree, Random Forest, and Linear SVC on unmodified and modified life-insurance datasets using accuracy, precision, recall, and F1-score.
    \item Identified Random Forest as the strongest classifier with 53.79\% accuracy on the unmodified dataset and 52.80\% on the modified dataset.
\end{itemize}

% ---------- Education ----------
\section{Education}

\educationHeading
{Blekinge Institute of Technology}
{Karlskrona, Sweden}
{Master of Science in Computer Science \textit{-- 50\% tuition scholarship based on academic merit}}
{Aug 2023 -- May 2026}

\vspace{3pt}

\educationHeading
{Jawaharlal Nehru Technological University Kakinada}
{Kakinada, India}
{Bachelor of Technology in Computer Science and Engineering}
{Aug 2019 -- Jun 2023}

% ---------- Certifications ----------
\section{Selected Certifications}
\begin{tabularx}{\textwidth}{X r}
\textbf{Machine Learning} & Coursera \\
\textbf{Programming, Data Structures and Algorithms using Python} & NPTEL \\
\textbf{Python (Basic) Skill Assessment} & HackerRank \\
\textbf{Machine Learning to Deep Learning: Remote Sensing Data Classification} & IIRS / ISRO \\
\end{tabularx}

% ---------- Honors & Achievements ----------
\section{Honors \& Achievements}
\begin{itemize}
    \item Secured 2nd rank at college level in the Scaler Edge Apex 2021 Hackathon.
    \item Earned 6-star HackerRank recognition in Problem Solving and Python Programming.
\end{itemize}

\end{document}
